\newpage
\begin{center}
	\large\textbf{Минобрнауки России}
	
	\large\textbf{Юго-Западный государственный университет}
	\vskip 1em
	\normalsize{Кафедра программной инженерии}
	\vskip 1em
	\ifВКР{
		\begin{flushright}
			\begin{tabular}{p{.4\textwidth}}
				\centrow УТВЕРЖДАЮ: \\
				\centrow Заведующий кафедрой \\
				\hrulefill \\
				\setarstrut{\footnotesize}
				\centrow\footnotesize{(подпись, инициалы, фамилия)}\\
				\restorearstrut
				«\underline{\hspace{1cm}}»
				\underline{\hspace{3cm}}
				20\underline{\hspace{1cm}} г.\\
			\end{tabular}
		\end{flushright}
	}\fi
\end{center}
\vspace{1em}
\begin{center}
	\large
	\ifВКР{
		ЗАДАНИЕ НА ВЫПУСКНУЮ КВАЛИФИКАЦИОННУЮ РАБОТУ
		ПО ПРОГРАММЕ БАКАЛАВРИАТА}
	\else
	ЗАДАНИЕ НА КУРСОВУЮ РАБОТУ (ПРОЕКТ)
	\fi
	\normalsize
\end{center}
\vspace{1em}
{\parindent0pt
	Студента \АвторРод, шифр\ \Шифр, группа \Группа
	
	1. Тема «\Тема\ \ТемаВтораяСтрока»
	\ifВКР{
		утверждена приказом ректора ЮЗГУ от \ДатаПриказа\ № \НомерПриказа
	}\fi.
	
	2. Срок предоставления работы к защите \СрокПредоставления
	
	3. Исходные данные для создания программной системы:
	
	3.1. Перечень решаемых задач:}

\renewcommand\labelenumi{\theenumi)}

\begin{enumerate}
	\item провести анализ деятельности спортивной организации и выявить бизнес-процессы, подлежащие автоматизации;
	\item спроектировать структуру информационной системы на платформе 1С:Предприятие 8.3;
	\item разработать и реализовать документы, обеспечивающие планирование тренировочного процесса, организацию и проведение соревнований, оформление спортивных разрядов, складской учет спортивного инвентаря;
	\item реализовать алгоритм заполнения планировщика расписания;
	\item реализовать разграничение прав доступа (роли «Администратор» и «Тренер») и провести функциональное тестирование системы.
\end{enumerate}

{\parindent0pt
	3.2. Входные данные и требуемые результаты для программы:}

\begin{enumerate}
	\item Входными данными для программной системы являются: данные о спортсменах и тренерах; данные о видах спорта, дисциплинах и разрядах; данные об учебных группах; временные интервалы для составления расписания; сведения о помещениях; данные о соревнованиях и их результатах; данные о спортивном инвентаре и его движении; файлы Excel с результатами соревнований.
	\item Выходными данными для программной системы являются: расписание занятий на месяц с детализацией по дням недели; отчеты по расписанию (общее, по группам, по тренерам); отчет по спортсменам с группировкой по видам спорта; отчет по результатам соревнований; печатные формы приказов о присвоении разрядов; текущие остатки спортивного инвентаря; уведомления участникам соревнований по электронной почте.
\end{enumerate}

{\parindent0pt
	
	4. Содержание работы (по разделам):
	
	4.1. Введение.
	
	4.2. Анализ предметной области.
	
	4.3. Техническое задание: основание для разработки, назначение разработки,
	требования к программной системе, моделирование вариантов использования, требования к оформлению документации.
	
	4.4. Технический проект: общая архитектура системы, структура конфигурации, структура справочников, структура документов, структура регистров, обоснование выбора платформы, разграничение прав доступа.
	
	4.5. Рабочий проект: спецификация программных модулей, реализация ключевых алгоритмов, функциональное тестирование системы.
	
	4.6. Заключение.
	
	4.7. Список использованных источников.
	
	5. Перечень графического материала:
	
	\списокПлакатов
	
	\vskip 2em
	\begin{tabular}{p{6.8cm}C{3.8cm}C{4.8cm}}
		Руководитель \ifВКР{ВКР}\else работы (проекта) \fi & \lhrulefill{\fill} & \fillcenter\Руководитель\\
		\setarstrut{\footnotesize}
		& \footnotesize{(подпись, дата)} & \footnotesize{(инициалы, фамилия)}\\
		\restorearstrut
		Задание принял к исполнению & \lhrulefill{\fill} & \fillcenter\Автор\\
		\setarstrut{\footnotesize}
		& \footnotesize{(подпись, дата)} & \footnotesize{(инициалы, фамилия)}\\
		\restorearstrut
	\end{tabular}
}

\renewcommand\labelenumi{\theenumi.}